% ============================================
% Chapter 6: Project Conclusion
% ============================================
\chapter{Project Conclusion}
\label{ch:conclusion}

This chapter concludes the thesis from a \textbf{project-level perspective}. The focus is not only on individual feature improvements, but on how RSHub-agent is built and validated as an end-to-end engineering system for remote sensing task assistance.

\section{Overall Project Outcomes}
Compared with the summer-project baseline, this iteration establishes a clearer full-stack path from user request to platform execution:
\begin{itemize}
	\item frontend/backend integration for conversational and streaming interaction
	\item ReAct-based orchestration with explicit HITL state control
	\item modular skill and tool systems for scalable capability extension
	\item engineering-side automation for synchronization and maintenance support
\end{itemize}

These outcomes align with the objective defined in Chapter~1: presenting RSHub-agent as a complete project implementation rather than a set of isolated patches.

\section{Core System Capabilities}

\subsection{Integrated Project Architecture}
In the final implementation, RSHub-agent is organized as an integrated project architecture rather than isolated scripts. The system connects frontend interaction, backend orchestration, skill and tool registries, and external platform services through explicit interfaces.

This architecture-level integration enables a stable end-to-end path from user intent parsing to remote task execution and result return, which is the structural basis for the workflow and engineering mechanisms described below.

\subsection{Reliable Workflow Control}
RSHub-agent enforces a controlled submission path through HITL gating. The operational flow remains explicit:
\begin{enumerate}
	\item user provides task intent and parameters
	\item agent generates executable plan
	\item HITL requests \texttt{CONFIRM}/\texttt{CANCEL}
	\item submission proceeds only after confirmation
	\item credit is deducted only on successful submission
\end{enumerate}

The billing rule is intentionally simple and transparent:
\begin{itemize}
	\item successful confirmed submission: deduct 1 credit
	\item cancel/timeout/non-matching input: deduct 0 credit
\end{itemize}

\subsection{Scalable Engineering Mechanisms}
The GitSync component provides repository-state-aware synchronization behavior, so the system does not rely on a single forced sync strategy:
\begin{itemize}
    \item behind only: fast-forward update (\texttt{merge --ff-only})
    \item ahead only: push requires explicit confirmation
    \item diverged: merge/rebase/force-with-lease/skip is selected explicitly
    \item unresolved conflicts: workflow exits with clear warning
\end{itemize}

\begin{lstlisting}[language=bash, basicstyle=\ttfamily\footnotesize]
git fetch origin main
git rev-list --left-right --count HEAD...origin/main
git merge --ff-only origin/main
\end{lstlisting}

This behavior improves operational reproducibility and reduces branch-state handling mistakes during routine maintenance.

In addition, the deployed skill architecture uses tiered loading:
\begin{itemize}
    \item \textbf{Layer 1 (always loaded)}: compact policy/index skills
    \item \textbf{Layer 2 (on-demand)}: detailed domain skills loaded per turn
\end{itemize}

In practical usage, this structure lowers redundant prompt payload and improves instruction relevance in multi-turn, task-oriented sessions. Quantitative confirmation remains in Chapter~\ref{ch:testing-report}.

\section{Project Methodology and Engineering Practices}
\label{sec:optimization-practices}

The engineering methodology in this project follows a repeatable pattern:
identify recurring operational risks, define explicit behavior rules, implement them in modular components, and validate through scenario-based testing.

\subsection{Practice A: Workflow Safety as a Primary Constraint}
\begin{itemize}
	\item Introduced explicit \texttt{CONFIRM}/\texttt{CANCEL} gating before task submission
	\item Applied charge only on successful submission (1 credit)
	\item Added strict keyword matching to prevent accidental submit triggers
\end{itemize}

\subsection{Practice B: Release Reliability through Automation}
\begin{itemize}
	\item Added repository-state checks before taking sync actions
	\item Restricted dangerous operations behind explicit user confirmation
	\item Kept conflict paths fail-fast with visible warnings
\end{itemize}

\subsection{Practice C: Context Efficiency through Skill Structuring}
\begin{itemize}
	\item Replaced monolithic system prompts with layer-1 + selective layer-2 skill loading
	\item Bound skill injection to turn intent (task/literature/HITL state)
	\item Reduced unnecessary context payload in non-relevant turns
\end{itemize}

\section{Future Work}
\label{sec:future-work}

Future work focuses on two major gaps in the current project iteration.
First, synchronization remains command-driven and human-triggered; a continuous maintenance loop (\textit{diff} detection, impact analysis, patch proposal, regression checks) is still incomplete.
Second, evaluation currently emphasizes workflow correctness more than long-run operational statistics under sustained workloads.

Another direction is an administrator-facing engineering assistant for repository hygiene and configuration consistency checks.
The goal is not generic coding help, but project-specific diagnostics: identifying risky configuration drift, incomplete test evidence, and repetitive maintenance tasks that can be standardized.

\section{Overall Summary}

This thesis presents RSHub-agent as an integrated remote sensing agent project spanning architecture, workflow control, skill organization, engineering automation, and software validation.
The current system is still an evolving engineering prototype, but its key behaviors and interfaces are now clearer, more testable, and more maintainable than in the baseline stage.
These results provide a practical foundation for scaling RSHub-agent toward more stable and extensible real-world deployment.
